% --- Chapter: System Programming with POSIX ---
\section{System Programming with POSIX}

CRISP provides full POSIX system call integration via an opt-in module. Unlike other scripting languages where system calls pollute the global namespace, CRISP loads POSIX on demand and namespaces all functions under a \texttt{posix} hash.

\subsection{Loading the POSIX Module}

\begin{tcolorbox}[colback=codebg, colframe=darkpurple, colbacktitle=darkpurple, title=\textbf{\textcolor{codegold}{Loading POSIX}}, fonttitle=\bfseries, sharp corners]
\begin{lstlisting}
# Load the POSIX module (opt-in, namespaced)
use posix;

# POSIX functions live under the posix hash
say posix["getpid"]();  # prints current PID

# Console read() is NOT shadowed!
let input = readline("> ");
say input;
\end{lstlisting}
\end{tcolorbox}

\subsection{Process Management}

\begin{tcolorbox}[colback=codebg, colframe=darkpurple, colbacktitle=darkpurple, title=\textbf{\textcolor{codegold}{Process Management}}, fonttitle=\bfseries, sharp corners]
\begin{lstlisting}
use posix;

# Fork a child process
let pid = posix["fork"]();
if pid == 0 {
    # Child process
    say "Child PID: ", posix["getpid"]();
    posix["execvp"]("ls", "-la");
} else {
    # Parent process
    say "Parent waiting for child: ", pid;
    let status = posix["waitpid"](pid, 0);
    say "Child exited with: ", status;
}

# Get process IDs
say "PID: ", posix["getpid"]();
say "PPID: ", posix["getppid"]();
\end{lstlisting}
\end{tcolorbox}

\subsection{File Operations (POSIX)}

\begin{tcolorbox}[colback=codebg, colframe=darkpurple, colbacktitle=darkpurple, title=\textbf{\textcolor{codegold}{POSIX File Operations}}, fonttitle=\bfseries, sharp corners]
\begin{lstlisting}
use posix;

# Low-level file descriptor operations
let fd = posix["open"]("/tmp/test.txt",
    posix["O_WRONLY"] | posix["O_CREAT"], 0o644);
posix["write"](fd, "Hello POSIX!\n", 13);
posix["close"](fd);

# Write to file
let fd = posix["open"]("data.txt",
    posix["O_WRONLY"] | posix["O_CREAT"], 0o644);
posix["write"](fd, "Hello POSIX!\n", 13);
posix["close"](fd);

# Seek in file
let fd2 = posix["open"]("data.bin", posix["O_RDONLY"]);
posix["lseek"](fd2, 0, posix["SEEK_END"]);
let size = posix["lseek"](fd2, 0, posix["SEEK_CUR"]);
\end{lstlisting}
\end{tcolorbox}

\subsection{Signals}

\begin{tcolorbox}[colback=codebg, colframe=darkpurple, colbacktitle=darkpurple, title=\textbf{\textcolor{codegold}{Signals}}, fonttitle=\bfseries, sharp corners]
\begin{lstlisting}
use posix;

# Send signal to a process
posix["kill"](pid, posix["SIGTERM"]);

# Set an alarm
posix["alarm"](5);  # SIGALRM in 5 seconds
\end{lstlisting}
\end{tcolorbox}

\subsection{Users and Environment}

\begin{tcolorbox}[colback=codebg, colframe=darkpurple, colbacktitle=darkpurple, title=\textbf{\textcolor{codegold}{Users and Environment}}, fonttitle=\bfseries, sharp corners]
\begin{lstlisting}
use posix;

# User/group IDs
say "UID: ", posix["getuid"]();
say "EUID: ", posix["geteuid"]();
say "GID: ", posix["getgid"]();

# Environment variables
let path = posix["getenv"]("PATH");
say "PATH: ", path;

posix["setenv"]("CRISP_TEST", "hello");
say posix["getenv"]("CRISP_TEST");  # hello

# Get all environment variables as hash
let env = posix["environ"]();
say env;
\end{lstlisting}
\end{tcolorbox}

\subsection{POSIX Function Reference}

\begin{table}[H]
\centering
\caption{\textcolor{darkpurple}{POSIX Functions}}
\vspace{0.3em}
\begin{tabular}{@{}lllp{3.2cm}@{}}
\toprule
\rowcolor{softvioletdark}
\textbf{\textcolor{white}{Category}} & \textbf{\textcolor{white}{Function}} & \textbf{\textcolor{white}{Signature}} & \textbf{\textcolor{white}{Description}} \\
\midrule
\rowcolor{lightlavender}
\multicolumn{4}{@{}l@{}}{\textbf{\textcolor{darkpurple}{Processes}}} \\
\rowcolor{softviolet!30}
& \texttt{fork} & \texttt{posix["fork"]() → Int} & Create child process \\
& \texttt{execvp} & \texttt{posix["execvp"](file, args...)} & Execute program \\
& \texttt{waitpid} & \texttt{posix["waitpid"](pid, opts) → Int} & Wait for child \\
& \texttt{spawn} & \texttt{posix["spawn"](cmd, args...)} & Spawn process \\
& \texttt{getpid} & \texttt{posix["getpid"]() → Int} & Current PID \\
& \texttt{getppid} & \texttt{posix["getppid"]() → Int} & Parent PID \\
\rowcolor{lightlavender}
\multicolumn{4}{@{}l@{}}{\textbf{\textcolor{darkpurple}{Signals}}} \\
\rowcolor{softviolet!30}
& \texttt{kill} & \texttt{posix["kill"](pid, sig)} & Send signal \\
& \texttt{alarm} & \texttt{posix["alarm"](secs)} & Schedule alarm \\
\rowcolor{lightlavender}
\multicolumn{4}{@{}l@{}}{\textbf{\textcolor{darkpurple}{Users}}} \\
\rowcolor{softviolet!30}
& \texttt{getuid} & \texttt{posix["getuid"]() → Int} & Real UID \\
& \texttt{geteuid} & \texttt{posix["geteuid"]() → Int} & Effective UID \\
& \texttt{getgid} & \texttt{posix["getgid"]() → Int} & Real GID \\
& \texttt{getegid} & \texttt{posix["getegid"]() → Int} & Effective GID \\
\rowcolor{lightlavender}
\multicolumn{4}{@{}l@{}}{\textbf{\textcolor{darkpurple}{Files}}} \\
\rowcolor{softviolet!30}
& \texttt{open} & \texttt{posix["open"](path, flags, mode?) → Int} & Open file \\
& \texttt{close} & \texttt{posix["close"](fd)} & Close file \\
& \texttt{read} & \texttt{posix["read"](fd, count) → Str} & Read bytes \\
& \texttt{write} & \texttt{posix["write"](fd, data, len) → Int} & Write bytes \\
& \texttt{lseek} & \texttt{posix["lseek"](fd, offset, whence) → Int} & Seek position \\
\rowcolor{lightlavender}
\multicolumn{4}{@{}l@{}}{\textbf{\textcolor{darkpurple}{Env}}} \\
\rowcolor{softviolet!30}
& \texttt{getenv} & \texttt{posix["getenv"](name) → Str} & Get env var \\
& \texttt{setenv} & \texttt{posix["setenv"](name, value)} & Set env var \\
& \texttt{environ} & \texttt{posix["environ"]() → Hash} & All env vars \\
\bottomrule
\end{tabular}
\end{table}

\begin{tcolorbox}[colback=lightlavender, colframe=softvioletdark, title=\textbf{\textcolor{darkpurple}{Design Note: Namespaced POSIX}}, fonttitle=\bfseries, sharp corners]
CRISP's decision to namespace POSIX under the \texttt{posix} hash solves the classic scripting language problem of namespace pollution. Languages like Perl and PHP expose hundreds of global functions, creating naming conflicts with user code. In CRISP:

\begin{itemize}
    \item \textbf{Opt-in:} POSIX is only loaded when \texttt{use posix;} is explicitly invoked
    \item \textbf{Namespaced:} All functions live under \texttt{posix["fn"]}, never in the global scope
    \item \textbf{No shadowing:} Console \texttt{read()} and POSIX \texttt{posix["read"]()} coexist without conflict
    \item \textbf{Discoverable:} \texttt{posix} is a plain Hash — iterate it with \texttt{for} or inspect with \texttt{keys()}
\end{itemize}

This pattern will extend to future modules: \texttt{use network;} creates a \texttt{network} hash, \texttt{use json;} creates a \texttt{json} hash, etc.
\end{tcolorbox}
